\subsection{The second six-generated family}
\label{sec:family-two}

In this subsection, we treat the second family \eqref{eq:family-two} in the
Introduction.  Put
$$
\begin{aligned}
 I={}&(g_1=y^2-x^{2b+1},\ g_2=x^{2b+2},\ g_3=x^{b+1}y,\
        g_4=x^rz,\ g_5=yz,\ g_6=z^2),\\
 &\hspace{35mm} b\geq1,\qquad 1\leq r\leq b+1.
\end{aligned}
$$

\subsubsection{Cohen--Macaulayness}

\begin{proposition}
\label{prop:family-two-reductions}
Define $Q\subseteq I$ by
$$
 Q=
 \begin{cases}
  (g_1,g_3+g_6,g_5),&r=b+1,\\
  (g_1,g_3+g_6,g_4),&1\leq r\leq b.
 \end{cases}
$$
Then the stronger equality $I^2=QI$ holds.  Consequently, $\calR(I)$ is Cohen--Macaulay by
Proposition~\ref{prop:local-reduction-criterion}.
\end{proposition}

\begin{proof}
We consider the two cases in the definition of $Q$.

\medskip

\noindent
\underline{$r=b+1$}: Since $Q=(g_1,g_3+g_6,g_5)$, we have
$I=Q+(g_2,g_4,g_6)$.  Hence it is enough to treat the six products
$g_ig_j$ with $i,j\in\{2,4,6\}$.
\begin{itemize}
 \item $g_2^2=xg_3(g_3+g_6)-xg_4g_5-xg_1g_2\in QI$.
 \item $g_2g_4=xg_3g_5-xg_1g_4\in QI$.
 \item $g_2g_6=xg_5^2-xg_1g_6\in QI$.
 \item $g_4^2=g_2g_6\in QI$.
 \item $g_4g_6=(g_3+g_6)g_4-g_2g_5\in QI$.
 \item $g_6^2=(g_3+g_6)g_6-g_4g_5\in QI$.
\end{itemize}
Thus $I^2=QI$.

\medskip

\noindent
\underline{$1\leq r\leq b$}: Since $Q=(g_1,g_3+g_6,g_4)$, we have
$I=Q+(g_2,g_5,g_6)$.  The six remaining products are treated as follows:
\begin{itemize}
 \item
 $g_2^2=xg_3(g_3+g_6)-x^{b+2-r}g_4g_5-xg_1g_2\in QI$.
 \item $g_2g_5=x^{b+1-r}g_3g_4\in QI$.
 \item $g_2g_6=x^{2b+2-2r}g_4^2\in QI$.
 \item $g_5^2=g_1g_6+x^{2b+1-2r}g_4^2\in QI$.
 \item
 $g_5g_6=(g_3+g_6)g_5-x^{b+1-r}g_1g_4-x^{b-r}g_2g_4\in QI$.
 \item $g_6^2=(g_3+g_6)g_6-x^{b+1-r}g_4g_5\in QI$.
\end{itemize}
Thus $I^2=QI$ also in this case.  In particular, unlike the calculation for
the first family, these choices of $Q$ are reductions of $I$ already in $S$.
\end{proof}

\subsubsection{Normality}

We next prove that $I$ is normal.  Recall from
Subsection~\ref{subsec:symmetric-multiplicity-four} that we have chosen the
grading given by $\deg x=4$, $\deg y=4b+2$, and
$\deg z=8b+7-4r$.  For this grading, the least relation degree is
$\ell=8b+4$.  The ideal $I$ itself is homogeneous for any positive choice of
$\deg z$.

Let $f=y^2$ and $J=I+(f)$.  We have
$$
 J=(y^2,x^{2b+1},x^{b+1}y,x^rz,yz,z^2).
$$
As observed in Section~\ref{sec:target-ideals}, $J=I_\ell$ and hence $J$
is integrally closed.  Since $J$ is a monomial ideal with six minimal
generators, it is normal by \cite{AtakaMatsuoka2026}.

\begin{lemma}\label{lem:common_secondfamily}
The following assertions hold.
\begin{enumerate}
 \item $\m f\subseteq I$.
 \item $fg_2,fg_6\in I^2$.
 \item If $r=b+1$, then $fg_4\in I^2$.
 \item If $1\leq r\leq b$, then $fg_5\in I^2$.
\end{enumerate}
\end{lemma}

\begin{proof}
The first assertion follows from
$$
 xf=xg_1+g_2,\qquad
 yf=yg_1+x^bg_3,\qquad
 zf=zg_1+x^{2b+1-r}g_4.
$$
The remaining assertions follow from
$$
 fg_2=g_3^2,\qquad fg_6=g_5^2,
$$
together with
$$
 fg_4=g_3g_5\quad(r=b+1),\qquad
 fg_5=g_1g_5+x^{b-r}g_3g_4\quad(1\leq r\leq b).
$$
\end{proof}

We apply the strategy described in Subsection~\ref{subsec:normality}.
Suppose that $I^s$ is not integrally closed for some positive integer $s$.
Choose $\xi$, $m$, and $\omega$ as in that subsection, so that
$$
 \xi=\xi_0+f^m\omega,
$$
where
$\xi_0\in I^s+fI^{s-1}+\cdots+f^{m-1}I^{s-m+1}$ and
$\omega\in I^{s-m}$.  By Lemma~\ref{lem:common_secondfamily}, we may write
$$
 \omega=\sum_{\lambda\in\Lambda}c_\lambda g^\lambda,
 \qquad c_\lambda\in k,
$$
where
$$
 \Lambda=
 \begin{cases}
 \{\lambda\in\Lambda^{s-m}_{\delta-m\ell}
       \mid\lambda_2=\lambda_4=\lambda_6=0\},&r=b+1,\\
 \{\lambda\in\Lambda^{s-m}_{\delta-m\ell}
       \mid\lambda_2=\lambda_5=\lambda_6=0\},&1\leq r\leq b.
 \end{cases}
$$

Put $\nu=4b+3$, and take $\varepsilon_1,\varepsilon_2\in k^\times$.  If
$\ch k\neq2$, define a homomorphism
$\psi_{\varepsilon_1,\varepsilon_2}\colon S\to k[[t]]$ by
$$
 x\longmapsto t^2,\qquad
 y\longmapsto t^{2b+1}+\varepsilon_1t^{2b+2}.
$$
If $\ch k=2$, use instead
$$
 x\longmapsto t^2(1+\varepsilon_1t),\qquad
 y\longmapsto t^{2b+1}.
$$
In either case, complete the definition by setting
$$
 z\longmapsto
 \begin{cases}
  \varepsilon_2t^{2b+2},&r=b+1,\\
  \varepsilon_2t^{\nu-2r},&1\leq r\leq b.
 \end{cases}
$$
Let $\kappa=2$ if $\ch k\neq2$, and let $\kappa=1$ if
$\ch k=2$.  For the elements that occur in $\omega$, the
leading terms are as follows:
\[
\begin{array}{c|c|c}
 \text{Element}&\psi_{\varepsilon_1,\varepsilon_2}&
 \ord_t\psi_{\varepsilon_1,\varepsilon_2}\\ \hline
 f=y^2&t^{\nu-1}+\text{(higher-order terms)}&\nu-1\\
 g_1&\kappa\varepsilon_1t^\nu+\text{(higher-order terms)}&\nu\\
 g_3&t^\nu+\text{(higher-order terms)}&\nu\\
 g_5\quad(r=b+1)&\varepsilon_2t^\nu+\text{(higher-order terms)}&\nu\\
 g_4\quad(1\leq r\leq b)&\varepsilon_2t^\nu+\text{(higher-order terms)}&\nu
\end{array}
\]
Moreover, $\psi_{\varepsilon_1,\varepsilon_2}(I)k[[t]]\subseteq(t^\nu)$.

Suppose first that $r=b+1$. The omitted generators $g_2,g_4,g_6$
have orders $\nu+1,\nu+1,\nu+1$, respectively. Set
\[
 G=\{i\mid\ord_t\psi_{\varepsilon_1,\varepsilon_2}(g_i)=\nu\}=\{1,3,5\}.
\]
Every $\lambda\in\Lambda$ has support in $G$.
In Lemma~\ref{lem:coefficient-separation}, the parameter exponent map is
$\lambda\mapsto(\lambda_1,\lambda_5)$.

The parameter exponents determine $\lambda_1$ and $\lambda_5$, respectively, and the
number-of-factors equation
$$
 \lambda_1+\lambda_3+\lambda_5=s-m
$$
then determines $\lambda_3$.  Thus distinct elements of $\Lambda$ give
distinct monomials in $\varepsilon_1$ and $\varepsilon_2$, and hence
$c_\lambda=0$ for every
$\lambda\in\Lambda$.

If $1\leq r\leq b$, the omitted generators $g_2,g_5,g_6$ have orders
$\nu+1$, $\nu+2b+1-2r$, and $\nu+4b+3-4r$, respectively.
Set
\[
 G=\{i\mid\ord_t\psi_{\varepsilon_1,\varepsilon_2}(g_i)=\nu\}=\{1,3,4\}.
\]
Every $\lambda\in\Lambda$ has support in $G$. Apply Lemma~\ref{lem:coefficient-separation}
with parameter exponent map $\lambda\mapsto(\lambda_1,\lambda_4)$.

Here the exponents of $\varepsilon_1$ and $\varepsilon_2$ determine
$\lambda_1$ and $\lambda_4$, while
$$
 \lambda_1+\lambda_3+\lambda_4=s-m
$$
determines $\lambda_3$.  Therefore every coefficient $c_\lambda$ is again
zero.

In both cases, $\omega=0$, contrary to the minimality of $m$.  It follows
that $I^s$ is integrally closed for every $s\geq1$, and hence $I$ is normal.

Combining this with Proposition~\ref{prop:family-two-reductions}, we obtain
the following.

\begin{theorem}
\label{thm:family-two}
Let $I$ be an ideal in the family~\eqref{eq:family-two}.  Then the Rees
algebra $\calR(I)$ is a Cohen--Macaulay normal domain.
\end{theorem}
